% Chapter 6: inference methodology for the late-stopped CG--FPLS path.
\chapter{Inference Methodology}
\label{chap:inference-methodology}

Chapter~\ref{chap:estimation-results} showed that stopping CG--FPLS early can
improve estimation of the slope function.  Hypothesis testing places a
different demand on the method.  If the iterations stop too soon, the error
remaining in the fitted equation can affect the test result.  The procedure
of \textcite{babii_carrasco_tsafack_2026} therefore continues until the fitted
empirical moment $\widehat K\widehat\beta_m$ is sufficiently close to
$\widehat r$: their difference must become negligible relative to the
root-$n$ sampling fluctuation.  This chapter explains why that requirement
matters and how it is assessed.

The analysis focuses on the moment-residual CG--FPLS method defined in
Section~\ref{sec:methods-cg}.  Its conclusions apply to this method and do
not automatically extend to raw FPLS, Arnoldi response-space FPLS, or FPCR
from Chapter~\ref{chap:estimation-results}.  The chapter introduces the
hypothesis and test statistic, explains the stopping requirement and choice
of critical value, and then considers power and confidence sets.  The final
section describes the simulation design used to assess these ideas in
Chapter~\ref{chap:inference-results}.

\section{Hypothesis and test statistic}
\label{sec:inference-target}

The starting point is the centred scalar-on-function model
\begin{equation}
    Y_i^c=\langle X_i^c,\beta\rangle+\varepsilon_i,
    \qquad
    \mathbb E(\varepsilon_i\mid X_i)=0,
    \label{eq:inference-model}
\end{equation}
and the normal equation $K\beta=r$.  The empirical moments
$\widehat K$ and $\widehat r$ are defined in
\eqref{eq:methods-empirical-moments}, with their discrete versions given in
\eqref{eq:methods-grid-moments}.  For a specified candidate slope $b$, the
question is whether the true slope equals this function:
\begin{equation}
    H_0:\beta=b
    \qquad\text{against}\qquad
    H_1:\beta\neq b.
    \label{eq:inference-hypothesis}
\end{equation}
The hypothesis concerns the whole slope function.  It does not ask only
whether $\beta(s)$ takes a particular value at one point.  Choosing $b=0$
tests whether the functional predictor has any effect, although other
candidate slopes can also be tested.  To assess how often the test rejects
a true null hypothesis, the simulation uses the known generating slope as $b$.

There is an identification limit, as explained in
Section~\ref{sec:identification-illposedness}.  If $K$ is not injective, the
observations determine only $K\beta$.  They cannot distinguish two slopes
whose difference lies in $\operatorname{Null}(K)$.  Without further
restrictions, the hypothesis that can be tested is therefore
\begin{equation}
    H_0:K\beta=Kb.
    \label{eq:inference-moment-null}
\end{equation}
Restricting $\beta$ and $b$ to
$\mathcal H_{\mathrm{id}}=\operatorname{Null}(K)^\perp$ makes this hypothesis
equivalent to \eqref{eq:inference-hypothesis}.  References to testing the
slope below use this restriction.  A change lying entirely in
$\operatorname{Null}(K)$ cannot be detected by a test based on the functional
linear model.

Let $\widehat\beta_m^{\mathrm{CG}}$ be the estimate that minimises the
empirical moment residual in the $m$th Krylov space, as defined in
\eqref{eq:methods-cg-moment-objective}.  The test statistic after $m$
iterations is
\begin{equation}
    \mathcal T_{n,m}(b)
    =n\left\|
       \widehat K\bigl(\widehat\beta_m^{\mathrm{CG}}-b\bigr)
      \right\|^2.
    \label{eq:inference-finite-m-statistic}
\end{equation}
When the stopping index depends on sample size, write $m=m_n$ and use
$\mathcal T_n(b)=\mathcal T_{n,m_n}(b)$, as in
\eqref{eq:intro-test-statistic}.  The calligraphic symbol $\mathcal T_n$
denotes this scalar statistic; $T_n$ continues to denote the sampling
operator introduced in Chapter~\ref{chap:methods}.

The multiplication by $\widehat K$ has a practical purpose.  A test based
directly on $\|\widehat\beta_m-b\|$ would be affected by the instability of
estimating slope components associated with small covariance eigenvalues.
For $\widehat K(\widehat\beta_m-b)$, the sampling fluctuation has a root-$n$
limit under the moment conditions stated below.  This allows inference about
the identifiable part of the slope without requiring the unweighted slope
estimate itself to have such a limit.

\section{Why inference requires later stopping}
\label{sec:inference-late-stopping}

Two differences explain the role of stopping.  The first measures how well
the candidate slope $b$ satisfies the empirical equation; the second is the
error remaining after $m$ CG iterations:
\begin{equation}
    u_n(b)=\widehat r-\widehat Kb,
    \qquad
    e_m=\widehat r-\widehat K\widehat\beta_m^{\mathrm{CG}}.
    \label{eq:inference-two-residuals}
\end{equation}
They give the exact identity introduced in
\eqref{eq:inference-residual-decomposition-background}:
\begin{equation}
    \widehat K(\widehat\beta_m^{\mathrm{CG}}-b)
    =u_n(b)-e_m.
    \label{eq:inference-residual-decomposition}
\end{equation}
Under the null, $u_n(b)$ reflects sampling variation, while $e_m$ reflects
the error left by stopping the iterations.  A statistic based directly on
the empirical equation is
\begin{equation}
    \mathcal S_n(b)
    =n\|u_n(b)\|^2
    =n\|\widehat r-\widehat Kb\|^2.
    \label{eq:inference-direct-statistic}
\end{equation}
Expanding the squared norm in \eqref{eq:inference-residual-decomposition}
gives
\begin{equation}
    \mathcal T_{n,m}(b)-\mathcal S_n(b)
    =-2n\langle u_n(b),e_m\rangle+n\|e_m\|^2,
    \label{eq:inference-statistic-difference}
\end{equation}
and hence
\begin{equation}
 \left|\mathcal T_{n,m}(b)-\mathcal S_n(b)\right|
 \leq
 2\sqrt n\|u_n(b)\|\sqrt n\|e_m\|
 +n\|e_m\|^2.
 \label{eq:inference-statistic-bound}
\end{equation}

Under the null, $\sqrt n\|u_n(b)\|=O_{\mathbb P}(1)$.  The bound in
\eqref{eq:inference-statistic-bound} therefore shows that the two statistics
become equivalent if the fitting residual decreases faster than $n^{-1/2}$.
This gives the sufficient late-stopping condition
\begin{equation}
    \sqrt n\,
    \|\widehat r-\widehat K\widehat\beta_{m_n}^{\mathrm{CG}}\|
    =o_{\mathbb P}(1).
    \label{eq:inference-late-stopping-condition}
\end{equation}
When this condition holds,
$\mathcal T_{n,m_n}(b)=\mathcal S_n(b)+o_{\mathbb P}(1)$: the error caused by
stopping no longer affects the limiting null distribution.  The argument is
given in Proposition~\ref{prop:app-inference-decomposition} and
Theorem~\ref{thm:app-late-stopping}.

The condition in \eqref{eq:inference-late-stopping-condition} is stronger
than the estimation rule in
\eqref{eq:methods-cg-stopping-rule}.  That rule stops when $\|e_m\|$ falls
below a data-dependent constant multiple of $n^{-1/2}$.  Such a threshold
can provide useful regularisation, but it does not ensure that
$\sqrt n\|e_m\|$ tends to zero.  Additional iterations can therefore make
the unweighted slope estimate less accurate while making the test statistic
more accurate.  The two tasks place different demands on the fitted moment.

The number of iterations needed depends on the covariance spectrum, the
sample and resulting Krylov spaces, and numerical precision.  Following the
paper's null-distribution experiment, the present null-law study and the
fixed-index stopping comparison use $m_{\mathrm{late}}=70$ at $n=100$ as a
conservative choice to be checked.  Power is studied using the direct-moment
statistic, while the confidence-set illustration uses $m_C=10$.  Neither
fixed iteration count replaces the requirement in
\eqref{eq:inference-late-stopping-condition}.

At the full empirical Krylov dimension, the residual is zero in exact
arithmetic whenever $\widehat r$ lies in
$\operatorname{Range}(\widehat K)$.  Reaching that dimension is unnecessary
for first-order inference if the residual is already small enough, and can
be undesirable for slope estimation.  The stopping study therefore examines
how the residual and test statistic change as the iterations continue.

\section{The null distribution and choice of critical value}
\label{sec:inference-calibration}

Suppose the observations are independent and identically distributed,
$\mathbb E\|X^c\|^4<\infty$, and
$\mathbb E(\varepsilon^2\mid X)\leq\sigma_{\max}^2<\infty$.  Under
$H_0:K\beta=Kb$, and with the known zero means used in the simulations, the
empirical null moment is exactly
\begin{equation}
    u_n(b)=\frac1n\sum_{i=1}^{n}\varepsilon_iX_i^c.
    \label{eq:inference-null-score-average}
\end{equation}
If sample means are used instead, the additional centring term is
asymptotically negligible under these moment conditions.  The covariance
operator of the random score $\varepsilon X^c$ is
\begin{equation}
    V
    =\mathbb E\!\left\{\varepsilon^2
       (X^c\otimes X^c)\right\},
    \qquad
    Vh=\mathbb E\!\left[\varepsilon^2
       \langle X^c,h\rangle X^c\right].
    \label{eq:inference-variance-operator}
\end{equation}
The Hilbert-space central limit theorem yields
\begin{equation}
    \sqrt n\,u_n(b)\rightsquigarrow G,
    \qquad G\sim\mathcal N_{\mathcal H}(0,V).
    \label{eq:inference-hilbert-clt}
\end{equation}
When the late-stopping condition holds, the test statistic has the same
limit as $\mathcal S_n(b)$.  Letting $\{\omega_j\}$ denote the positive
eigenvalues of $V$, the spectral representation of $G$ and Parseval's
identity give
\begin{equation}
    \mathcal T_n(b)
    \rightsquigarrow
    \|G\|^2
    =\sum_{j\geq1}\omega_jZ_j^2,
    \qquad Z_j\stackrel{\mathrm{iid}}{\sim}\mathcal N(0,1).
    \label{eq:inference-weighted-chi-square}
\end{equation}
The series is well defined because $V$ is trace class under the stated
moment conditions.  Its weights depend on the joint distribution of $X$ and
the regression error, so the limiting law contains unknown population
quantities.  This is why a critical value cannot simply be taken from an
ordinary chi-square distribution with known degrees of freedom.

When the errors are conditionally homoskedastic,
$\mathbb E(\varepsilon^2\mid X)=\sigma^2$, then
\begin{equation}
    V=\sigma^2K,
    \qquad
    \omega_j=\sigma^2\lambda_j,
    \label{eq:inference-homoskedastic-weights}
\end{equation}
where $\{\lambda_j\}$ are the positive eigenvalues of $K$.  Under
heteroskedasticity, $V$ and $K$ need not share eigenfunctions.  The
homoskedastic simplification is proved in
Proposition~\ref{prop:app-homoskedastic-variance}.

Let $q_{1-\alpha}$ denote the $(1-\alpha)$ quantile of the law in
\eqref{eq:inference-weighted-chi-square}.  The level-$\alpha$ rule is
\begin{equation}
    \varphi_n(b)
    =\mathbf 1\{\mathcal T_n(b)>q_{1-\alpha}\}.
    \label{eq:inference-rejection-rule}
\end{equation}
If the limiting law is non-degenerate, it is continuous at this quantile.
The probability of rejecting a true null hypothesis then approaches
$\alpha$, as stated in Corollary~\ref{cor:app-asymptotic-size}.  This is an
asymptotic result; it does not guarantee an exact rejection rate in a finite
sample.

The study distinguishes three ways of choosing critical values.

\paragraph{Using the known simulation model: oracle calibration.}
The simulation specifies the score variances of the continuous model and
the error variance.  These known values allow independent $Z_j$ to be drawn
and combined as
\begin{equation}
    Q^{\mathrm{oracle}}
    =\sum_{j=1}^{J}\sigma^2\lambda_jZ_j^2,
    \label{eq:inference-oracle-law}
\end{equation}
and the critical value is estimated from $50{,}000$ draws of this quantity.
The term \emph{oracle} indicates that the calculation uses population
information supplied by the simulation design.  Even with this information,
the test need not have exact size at finite $n$.

There is also a distinction between the continuous model and its numerical
grid.  The equally weighted grid includes both endpoints, so the basis
functions are only approximately orthonormal under the grid inner product.
To describe the resulting difference, let
$\Phi=(\phi_j(s_t))_{j,t}\in\mathbb R^{J\times T}$,
$G_T=T^{-1}\Phi\Phi^\top$, and
$\Lambda=\operatorname{diag}(\lambda_1,\ldots,\lambda_J)$.
Under the grid inner product, the non-zero covariance eigenvalues are those
of $\Lambda^{1/2}G_T\Lambda^{1/2}$.  Writing these eigenvalues as
$\nu_{T,j}$, with zeros included if necessary, gives the fixed-grid limit
\begin{equation}
    Q_T=\sum_{j=1}^{J}\sigma^2\nu_{T,j}Z_j^2.
    \label{eq:inference-grid-oracle-law}
\end{equation}
With $T=200$ held fixed, a larger sample does not remove the difference
between $\nu_{T,j}$ and $\lambda_j$.  For fixed $J$, a finer grid gives
$G_T\to I_J$ and recovers the continuous-model reference.  The reported
experiment uses \eqref{eq:inference-oracle-law}.  Any disagreement with that
reference can therefore reflect finite-sample effects, grid approximation,
and Monte Carlo error in the simulated distribution and critical value.
It cannot be attributed solely to finite-sample error relative to the exact
fixed-grid limit.

The oracle draws are generated independently for each model.  Models with
the same theoretical weights may consequently receive slightly different
estimated critical values.  These small differences arise from simulation
uncertainty and do not indicate different theoretical null distributions.

\paragraph{Using simulated null samples: empirical-null calibration.}
A second approach repeatedly generates samples under a known null
hypothesis and takes the required quantile of their test statistics.
This avoids the continuous-model weighted chi-square approximation,
although the estimated critical value still has Monte Carlo uncertainty.
The power study uses this approach separately for each model and sample
size.  It is available because the null model is known in the simulation;
an observed dataset alone does not provide independent samples from a known
null model.

\paragraph{Estimating critical values from observed data.}
In an application, the unknown features of the null distribution must be
estimated.  For a specified null, define the residual
$\widehat\varepsilon_i(b)=Y_i^c-\langle X_i^c,b\rangle$.  One route replaces
$V$ by the covariance estimator
\begin{equation}
    \widehat V
    =\frac1n\sum_{i=1}^{n}
      \widehat\varepsilon_i(b)^2
      (X_i^c\otimes X_i^c)
    \label{eq:inference-plugin-variance}
\end{equation}
and uses its estimated eigenvalues to simulate a weighted chi-square
distribution.  A non-parametric bootstrap offers another way to approximate
the null law.  Both approaches are discussed by
\textcite{babii_carrasco_tsafack_2026}.  A centred-multiplier method is also
available as an exploratory option in the accompanying implementation.
The reported null-size and confidence-set results, however, use only oracle
calibration.  Chapter~\ref{chap:inference-results} therefore provides no
empirical assessment of these methods for estimating critical values from
observed data.

\section{Alternatives and power}
\label{sec:inference-power}

Power is the probability that a test detects a departure from the null
hypothesis.  For a fixed alternative with $K(\beta-b)\neq0$, and with the
required moment convergence and sufficiently late stopping,
\begin{equation}
    \frac1n\mathcal T_n(b)
    \xrightarrow{\mathbb P}
    \|K(\beta-b)\|^2>0.
    \label{eq:inference-fixed-alternative}
\end{equation}
The statistic therefore grows without bound and the probability of
rejection approaches one.  The condition $K(\beta-b)\neq0$ excludes changes
in the slope that lie in the null space and are invisible in the data.

A more demanding question is whether the test can detect departures that
become smaller as the sample grows.  Consider
\begin{equation}
    H_{1,n}:\quad
    \beta_n=b+n^{-1/2}\Delta,
    \qquad \Delta\in\mathcal H.
    \label{eq:inference-local-alternative}
\end{equation}
Here the departure is in a fixed direction $\Delta$, with magnitude
decreasing at rate $n^{-1/2}$.  Under the same late-stopping condition,
\begin{equation}
    \mathcal T_n(b)
    \rightsquigarrow
    \|G+K\Delta\|^2.
    \label{eq:inference-local-limit}
\end{equation}
Writing $V\phi_j=\omega_j\phi_j$ makes the shift explicit:
\begin{align}
    \|G+K\Delta\|^2
    ={}&\sum_{j\geq1}\omega_jZ_j^2
       +2\sum_{j\geq1}\sqrt{\omega_j}Z_j
          \langle\phi_j,K\Delta\rangle
       +\|K\Delta\|^2.
    \label{eq:inference-local-expansion}
\end{align}
The change that enters the limiting statistic is $K\Delta$.  Thus the size
of a slope difference alone does not determine how easily it can be
detected.  Under homoskedasticity, writing $\Delta=\sum_j d_jv_j$ in the
eigenbasis of $K$ gives
\begin{equation}
    \|K\Delta\|^2
    =\sum_{j\geq1}\lambda_j^2d_j^2.
    \label{eq:inference-detectability}
\end{equation}
For slope changes with comparable coefficient norms, components associated
with larger covariance eigenvalues contribute more to this shift.
Departures in these directions are typically easier to detect than those
concentrated in weak directions.  More generally, power also depends on how
$K\Delta$ is oriented relative to the covariance operator $V$, so
$\|K\Delta\|$ alone does not determine it.  These local results concern
each fixed direction separately; they do not guarantee uniform power over
all unit-norm directions in an ill-posed problem.  The fixed and local
alternative results are proved in
Proposition~\ref{prop:app-inference-alternatives}.

The finite-sample power experiment uses the perturbation family
\begin{equation}
    \beta_\delta(s)=\beta(s)+\delta\Delta(s),
    \qquad
    \Delta(s)=s,
    \qquad
    \delta\in\{-1,-0.95,\ldots,0.95,1\}.
    \label{eq:inference-power-family}
\end{equation}
The candidate being tested remains the baseline $b=\beta$.  At $\delta=0$
the null is true; each fixed non-zero $\delta$ gives a fixed departure from
it at the chosen sample size.  To represent the shrinking alternatives in
\eqref{eq:inference-local-alternative}, $\delta$ would instead need to vary
with sample size as $\delta= c/\sqrt n$ for a fixed constant $c$.

The power study uses the direct statistic $\mathcal S_n(b)$ in
\eqref{eq:inference-direct-statistic}, evaluated exactly in the generating
score coordinates.  This keeps the large experiment manageable and measures
the finite-sample power of the statistic that sufficiently late-stopped
CG--FPLS approximates to first order.  The curves do not measure power after
a fixed $m=70$ iterations.  A separate stopping study checks the quality of
that approximation under the null.

For each model and $n\in\{100,200\}$, the power study uses an empirical null
quantile at the 5\% level.  Monte Carlo standard errors for the rejection
frequencies follow the binomial formula in \eqref{eq:simulation-mcse}.
They measure uncertainty conditional on the simulated critical values and
do not include the additional uncertainty from estimating an oracle or
empirical-null quantile.

\section{Confidence sets by test inversion}
\label{sec:inference-confidence}

A confidence set can be formed by collecting all candidate slopes that the
test does not reject.  For a level-$\alpha$ test on an identified parameter
space $\Theta\subseteq\mathcal H_{\mathrm{id}}$, this gives
\begin{equation}
    \mathcal C_{1-\alpha}
    =\left\{b\in\Theta:
      \mathcal T_n(b)\leq q_{1-\alpha}\right\}.
    \label{eq:inference-confidence-set}
\end{equation}
If the test has asymptotic size $\alpha$, then
$\mathbb P_\beta\{\beta\in\mathcal C_{1-\alpha}\}\to1-\alpha$; see
Corollary~\ref{cor:app-confidence-inversion}.  The identification restriction
matters here too.  Without it, adding a function in $\operatorname{Null}(K)$
does not change the test, so the set describes slopes that the data cannot
distinguish.  In a finite sample, $\widehat K$ also has finite rank.  An
unrestricted inverted set therefore extends without bound in directions in
$\operatorname{Null}(\widehat K)$.

To make the set easier to display, candidate functions are restricted to a
$p$-dimensional space:
\begin{equation}
    \mathcal H_p
    =\operatorname{span}\{h_1,\ldots,h_p\},
    \qquad
    b_c(s)=\sum_{j=1}^{p}c_jh_j(s),
    \label{eq:inference-truncated-space}
\end{equation}
where the illustration uses the first $p=5$ cosine functions from the
simulation basis.  Let $H$ be the $T\times p$ matrix whose $j$th column
contains $h_j$ on the numerical grid, and define
\begin{equation}
    A=\widehat K H,
    \qquad
    f=\widehat K\widehat\beta_{m_C}^{\mathrm{CG}},
    \label{eq:inference-confidence-matrices}
\end{equation}
and recall that the discrete squared $L^2$ norm is $T^{-1}\|\cdot\|_2^2$.
The accepted coefficient set is exactly
\begin{equation}
    C_{1-\alpha}^{(p)}
    =\left\{c\in\mathbb R^p:
       \frac nT\|f-Ac\|_2^2\leq q_{1-\alpha}
      \right\}.
    \label{eq:inference-coefficient-set}
\end{equation}

When $A$ has full column rank, let
\begin{equation}
    \widehat c=(A^\top A)^{-1}A^\top f,
    \qquad
    \rho^2=\frac{Tq_{1-\alpha}}n-
             \|f-A\widehat c\|_2^2.
    \label{eq:inference-ellipsoid-parameters}
\end{equation}
If $\rho^2<0$, no candidate in the restricted space is accepted and the set
is empty.  Otherwise, the accepted coefficients form the ellipsoid
\begin{equation}
    (c-\widehat c)^\top A^\top A(c-\widehat c)
    \leq\rho^2.
    \label{eq:inference-confidence-ellipsoid}
\end{equation}
\begin{samepage}
For the row vector $h(s)^\top=(h_1(s),\ldots,h_p(s))$, the smallest and largest
accepted function values at each point $s$ are
\begin{equation}
    h(s)^\top\widehat c
    \mathbin{\pm}
    \left\{
      \rho^2 h(s)^\top(A^\top A)^{-1}h(s)
    \right\}^{1/2}.
    \label{eq:inference-envelope}
\end{equation}
\end{samepage}
Proposition~\ref{prop:app-confidence-ellipsoid} derives
\eqref{eq:inference-confidence-ellipsoid}--\eqref{eq:inference-envelope} and
describes the cases where $A$ is rank deficient.  The analytic calculation
gives the exact envelope of the full ellipsoid in $\mathbb R^5$.  The
reference illustration instead searches over $20^5$ coefficient vectors in
$[0,4.5]^5$.  Both approaches invert the test, but the earlier search also
restricts the coefficients to a box and checks only finitely many points.
The displayed sets therefore need not coincide.

The confidence illustration uses $m_C=10$, following the separate choice in
the reference study; the null-law experiment uses
$m_{\mathrm{late}}=70$.  Ten components do not by themselves establish
\eqref{eq:inference-late-stopping-condition}.  Their adequacy must be assessed
through finite-sample coverage and residual diagnostics.  This choice is
therefore specific to the illustration, and its interpretation as a
confidence set remains conditional on a sufficiently small moment residual.

There is also a distinction between the full slope and its five-term
approximation.  The simulated true slopes have non-zero coefficients beyond
the fifth basis function, so they do not belong to the space used for the
display.  The study accordingly reports three different summaries:
\begin{enumerate}
    \item how often the test does not reject the full true simulated slope;
    \item how often the true slope's five-term projection is accepted within
          $\mathcal H_5$; and
    \item the lower and upper pointwise envelope of the accepted
          five-dimensional coefficient set.
\end{enumerate}
The shaded region in Chapter~\ref{chap:inference-results} combines the
pointwise medians of 5,000 lower and upper envelopes.  It shows a typical
range across many simulated samples.  It is not a 95\% confidence band from
one observed sample, and whether it contains the true curve does not measure
coverage.

The full-slope acceptance calculation uses the same samples as the null-size
experiment, but applies $m_C=10$ in place of $m_{\mathrm{late}}=70$.  Its
average can therefore be compared with the reported late-stopped size,
although the two need not add to one.  Acceptance and rejection are exact
complements only when the stopping index and critical value are the same.
Because the samples are shared, this comparison is not an independent check
of the size result.

\section{Implementation and diagnostic design}
\label{sec:inference-implementation}

The inference study uses the three models in
Sections~\ref{sec:simulation-dgp} and \ref{sec:simulation-models}: the same
100-term Gaussian score expansion, 200-point grid, Babii cosine convention,
and homoskedastic error standard deviation one.  Keeping these features
allows the inference results to be related to the estimation results.  A
separate random seed keeps the two simulation studies separate.

Table~\ref{tab:inference-design} gives the choices specific to inference.
The number of simulated samples applies to every model and, for power, to
every combination of sample size and perturbation value.  The size study examines
nominal levels of 1\%, 5\%, and 10\%; the other parts use 5\%.

\begin{table}[htbp]
    \centering
    \caption[Inference-study design]{Design of the CG--FPLS inference study.}
    \label{tab:inference-design}
    \small
    \begin{tabular}{@{}
        >{\raggedright\arraybackslash}p{2.75cm}
        >{\raggedright\arraybackslash}p{2.75cm}
        >{\raggedright\arraybackslash}p{3.1cm}
        >{\raggedright\arraybackslash}p{5.25cm}@{}}
        \toprule
        Study branch & Statistic and stopping choice & Sample and replications & Calibration and purpose \\
        \midrule
        Null law and size
        & $\mathcal T_{n,70}$
        & $n=100$; $R=5{,}000$
        & $50{,}000$-draw oracle weighted chi-square law at
          $\alpha\in\{0.01,0.05,0.10\}$; compare the finite-sample null law
          with the continuous-design reference. \\
        Stopping diagnostics
        & $\mathcal T_{n,m}$ and $\mathcal S_n$
        & $n=100$; $R=5{,}000$
        & Oracle 5\% critical value; compare
          $m\in\{1,2,3,5,10,20,40,70\}$, the estimation-oriented adaptive
          variant, and the maximum-rank fit. \\
        Power
        & Direct $\mathcal S_n$
        & $n\in\{100,200\}$; $R=5{,}000$ for every $\delta$
        & Empirical 5\% null quantile for each model and $n$; evaluate
          \eqref{eq:inference-power-family}. \\
        Inverted-set illustration
        & $\mathcal T_{n,10}$; analytic $p=5$ inversion
        & $n=100$; $R=5{,}000$
        & Oracle 5\% critical value; record full-slope non-rejection,
          five-term projected acceptance, rank, emptiness, and envelope width. \\
        \bottomrule
    \end{tabular}
\end{table}

The stopping comparison includes an adaptive rule intended for estimation.
It uses the discrepancy formula and constants from
Chapter~\ref{chap:methods}, with $\tau=1.01$ and
$\delta_{\mathrm{stop}}=0.1$.  Its FPCR moment-GCV variance pilot follows
\eqref{eq:simulation-moment-gcv}--\eqref{eq:simulation-cg-sigma-pilot}, but
uses at most 20 components, compared with 70 in the estimation experiment.
The comparison also uses the stabilised moment-residual path described below
in place of the estimation study's short recurrence.  These differences
matter: the smaller pilot cap can change the estimated variance, the
discrepancy threshold, and the selected stopping index.  This is therefore a
variant of the procedure in Chapter~\ref{chap:estimation-results}.

The notation $\delta_{\mathrm{stop}}$ distinguishes the tuning constant from
the power perturbation $\delta$ in \eqref{eq:inference-power-family}.  A
maximum-rank fit provides another benchmark by projecting $\widehat r$ onto
the numerical range of $\widehat K$.  For every stopping rule, the study
records
\begin{equation}
    \sqrt n\|e_m\|,
    \qquad
    |\mathcal T_{n,m}(b)-\mathcal S_n(b)|,
    \qquad
    \mathbf 1\{\mathcal T_{n,m}(b)>q_{0.95}\},
    \label{eq:inference-stopping-diagnostics}
\end{equation}
together with the selected or fixed index.  These quantities show how much
moment residual remains, how closely the fitted statistic matches the direct
statistic, and whether the test rejects.  They allow the effect of stopping
on null rejection rates to be assessed.

The moment-residual path is computed in a numerically stabilised form because
the inference study takes many more iterations than the estimation rule.
At each dimension, two-pass reorthogonalisation forms an orthonormal basis
$Q_m$ for $\widehat{\mathcal K}_m(\widehat K,\widehat r)$, and a QR solve
computes
\begin{equation}
    \widehat a_m
    \in\arg\min_{a\in\mathbb R^m}
       \|\widehat r-\widehat KQ_ma\|_2,
    \qquad
    \widehat\beta_m^{\mathrm{CG}}=Q_m\widehat a_m.
    \label{eq:inference-stable-engine}
\end{equation}
In exact arithmetic, this solves the same minimum-moment-residual problem at
each dimension as the original short recurrence.  The orthogonal basis
improves numerical stability and can change the path in finite precision.
It does not change the objective to that of Arnoldi FPLS in
Chapter~\ref{chap:methods}.  Both methods use an orthogonal Krylov basis, but
Arnoldi FPLS minimises response error, while
\eqref{eq:inference-stable-engine} minimises the CG--FPLS moment residual.
Propositions~\ref{prop:app-arnoldi-span} and
\ref{prop:app-response-moment-krylov} establish the relevant distinctions
between their spaces and objectives.

All random quantities are generated from the master seed 2025.  Separate
random streams are used for the stopping experiment and the alternative
power samples.  Some quantities are deliberately shared: the confidence
experiment uses the null-size samples, the $n=100$ empirical power
calibration uses the direct null statistics, and the same oracle draws serve
all relevant parts of the study.  These choices make the calculations
repeatable and clarify which comparisons share simulated information.

Interpreting the results first requires checking whether the remaining
moment residual is small enough for the fitted statistic to behave like the
direct statistic.  Chapter~\ref{chap:inference-results} therefore begins
with the stopping diagnostics, followed by null size, power, and
confidence-set behaviour.
